\scanpage{172}
\mspage{88}

dont on \deleted{\unsure{dé}}déduit $\widetilde U$ \deleted{comme $\operatorname{Ker}(B'\longrightarrow Z')$}, et $\widetilde B''$ comme .

\[
 \underbrace{\widetilde B''\simeq B'\times_{Z'}Z}_{\substack{\text{image inv. de}\\\text{l'ext $B'$ de $Z'$ par $\widetilde U$,}\\\text{par l'hom }Z\mathrel{\text{---}} Z'}},
 \qquad\widetilde U=\operatorname{Ker}(B'\mathrel{\unsure{\text{---}\bullet|\longrightarrow}} Z').\tag*{(48)}
\]

\deleted{Ceci dit,} \inserted{Il me semble} que c'est sous la forme (46), où $\widetilde B''$ figure déjà nommément, que la donnée est \inserted{cependant} la plus maniable, puisque c'est \deleted{\unsure{sur|sous}} $\widetilde B''$ qui intervient directement dans la donnée (42\deleted{\unsure{b}} c).

\deleted{\unsure{Il}} Pour expliciter enfin la donnée (42 d), en écrivant $B=B'\times Z$, et en notant que l'opération de $B$ sur $G_0$ doit être triviale sur $Z$, on voit que cette donnée se réduit à celle d'une op. de $B'$ sur $G_0$ ; \deleted{\unsure{i.e.} (33)} pour écrire que \deleted{cette} l'opération \inserted{de $B$ sur $\unsure{G|G_0}$} se réduit sur $\widetilde B''$ à l'opération tautologique déduite du plongement $\widetilde B''\subset G_0$, on note que l'opération en question sur $\widetilde B''$ est l'opération composée

\[ \widetilde B''\longrightarrow B'\longrightarrow\operatorname{Aut}_{\rm gr}(G_0) \tag*{(49)} \]

qui doit donc être l'opération tautologique.

En résumé, on trouve que les données (26) s'explicitent en termes de théorie des

% archive sheet number: 171
